% Lab 7: Lab test. Compile: latexmk -xelatex lab07.tex
\documentclass[10pt,aspectratio=169,t]{beamer}
\usetheme[lab]{upbminimal}

\course[ADMD]{Application Development for Mobile Devices}
\decklabel{Lab 7}
\title{Lab test}
\subtitle{Rules, format, task types}
\author{Dinu-Ștefan Rusu}
\date{Week 14}

\begin{document}

\titleframe

\begin{frame}[eyebrow=Today]{Lab objectives}
  \vskip 2mm
  \begin{grid}[2]
    \cell[\faIcon{stopwatch}]{Prove four skills, once}{Solve four tasks from a one-page spec in one hour: a screen, a Cubit, a network call, and a route.}
    \cell[\faIcon{graduation-cap}]{Worth 30 points}{7.5 points per task, graded at the end of the test.}
  \end{grid}
  \vskip 6mm
  \taskmeta{Time}{1 hour}{Submit}{To the lab instructor, when the time is up}
\end{frame}

\begin{frame}[eyebrow=Rules]{Rules}
  \begin{steps}
    \item You work alone. No help from anyone else, no copying a classmate's work.
    \item One hour, one sitting, starting when you receive the spec.
    \item You must be able to explain any line the grader points to, in your own words.
    \item Every task uses only what Labs 2 to 6 taught.
  \end{steps}
  \begin{takeaway}{The clock does not pause, and the hour is not extended.}
  \end{takeaway}
\end{frame}

\begin{frame}[eyebrow=Format]{Test format}
  \vskip 2mm
  \begin{grid}[2]
    \cell[\faIcon{file-alt}]{One page, handed at the start}{Four tasks. Nothing changes once the clock starts.}
    \cell[\faClock]{One hour, then you submit}{Work until time is up, then submit your answers to the lab instructor. No extensions.}
  \end{grid}
  \vskip 6mm
  \muted{Every spec has the same four kinds of task, described next.}
\end{frame}

\begin{frame}[eyebrow=Format]{Task types}
  \begin{grid}[2]
    \cell[\faThLarge]{A screen}{Built from the layout widgets: Row, Column, ListView, Card. No hardcoded pixel offsets.}
    \cell[\faToggleOn]{State in a Cubit}{An Equatable state class, \code{emit}, and a BlocBuilder. Not a raw setState.}
    \cell[\faWifi]{One network call}{A single HTTP call parsed into a model class with fromJson. Not a hardcoded map.}
    \cell[\faIcon{route}]{One navigation}{A go\_router route that carries data and pushes to a second screen.}
  \end{grid}
\end{frame}

\closingframe{Lab test}{Four tasks · one hour}

\end{document}
